\section{Lecture 3, Part 0: Discrete Random Variables}

\begin{frame}
\vfill
\centering
{\color{RLBlue}\large Lecture 3}

\vspace{0.3em}
{\color{RLBlue}\Large\bfseries Part 0}

\vspace{0.6em}
{\Large Discrete Random Variables}
\vfill
\end{frame}

\begin{frame}{Learning Goals}
\begin{itemize}
  \item Write a quantity that depends on chance as a \textbf{random variable}, and list what it can be with a probability mass function.
  \pause
  \item Compute the \textbf{expectation} and \textbf{variance} of a random variable, and say how each relates to the sample mean and sample variance of Lecture 1.
  \pause
  \item Predict what happens to both when the variable is shifted or rescaled.
\end{itemize}
\end{frame}

% ─── Random variables ───────────────────────────────────────────────────────
\begin{frame}{Why Do We Need a Random Variable?}
\begin{itemize}
  \item \textbf{Why:} so far every probability has been attached to an \emph{event} -- ``the die is even'', ``the student takes the bus''. But most questions are about a \emph{number}: how many, how long, how much.
  \pause
  \item Writing ``the number of customers who buy a pastry'' in full each time is unworkable. Give it a letter, $X$, and we can add it, average it, and compare it.
  \pause
  \item \textbf{Everyday:} roll two dice and let $X$ be the sum. Rolling is chance; $X$ is the number chance produced. Board games are decided by $X$, not by which die showed what.
  \pause
  \item \textbf{Discrete} means the possible values can be listed: $0, 1, 2, \dots$ Lecture 4 does the continuous case, where they cannot.
\end{itemize}
\end{frame}

\begin{frame}{Random Variables and the Probability Mass Function}
\begin{rldefinition}
A \textbf{random variable} $X$ attaches a number to every outcome. \emph{Two dice:} the outcome is the pair $(3,4)$, and $X$, their sum, is $7$.

\smallskip
Its \textbf{probability mass function} $p(x) = P(X = x)$ lists how likely each value is. Every $p(x) \ge 0$, and they add to $1$ because $X$ must take \emph{some} value.
\end{rldefinition}
\pause
\begin{itemize}
  \item \textbf{Running example -- a coffee shop.} Let $X$ be the number of pastries in one order:
  \[
    p(0) = 0.7, \qquad p(1) = 0.2, \qquad p(2) = 0.1,
    \qquad \textstyle\sum_x p(x) = 1 .
  \]
  \pause
  \item This is the frequency distribution of Lecture 1 with one change: those were \emph{counts} from a sample we measured, these are \emph{probabilities} for a customer we have not served yet.
\end{itemize}
\end{frame}

\begin{frame}{Real Event: Why 7 Wins at Board Games}
\begin{columns}[T]
\column{0.50\textwidth}
\small
\begin{itemize}
  \item Each bar is $p(x)$; the number above it is how many of the $36$ equally likely pairs give that sum.
  \pause
  \item There is one way to roll $2$ ($1{+}1$) but six ways to roll $7$ ($1{+}6$, $2{+}5$, \dots, $6{+}1$), so $p(7) = \tfrac{6}{36} = \tfrac16$ -- six times more likely than $2$.
  \pause
  \item \textbf{Why it is a real event:} Monopoly's orange properties sit six to nine squares past Jail -- the stretch a player leaving Jail is most likely to reach -- and backgammon openings are built on the same table.
\end{itemize}
\column{0.48\textwidth}
\centering
\includegraphics[width=\linewidth]{figures/discrete_rv/discrete_rv_two_dice_pmf.pdf}
\end{columns}
\end{frame}

% ─── Expectation ────────────────────────────────────────────────────────────
\begin{frame}{Why Do We Need Expectation?}
\begin{rldefinition}
The \textbf{expectation} $E[X]$ is the average value of $X$ in the long run, with each value weighted by how likely it is:
\[
E[X] = \sum_x x\,p(x) .
\]
It is also written $\mu$ -- the same Greek letter Lecture 1 used for a population mean, because that is exactly what it is.
\end{rldefinition}
\pause
\begin{itemize}
  \item \textbf{Coffee shop:} $E[X] = 0(0.7) + 1(0.2) + 2(0.1) = 0.4$ pastries per order -- so $100$ customers should buy about $40$ pastries.
  \pause
  \item \textbf{Note:} $0.4$ is not a possible order. An expectation is a long-run average, not a prediction of any single case -- like ``$1.8$ children per family''.
\end{itemize}
\end{frame}

\begin{frame}{Worked Example: Is the Raffle Ticket Worth It?}
\begin{itemize}
  \item A school sells $500$ raffle tickets at \$2 each. One ticket wins a \$400 bicycle.
  \pause
  \item Let $X$ be your \emph{net} gain. You either win the bicycle having paid \$2, or you just lose the \$2:
  \[
    p(398) = \tfrac{1}{500}, \qquad p(-2) = \tfrac{499}{500}.
  \]
  \pause
  \item $E[X] = 398 \cdot \tfrac{1}{500} - 2 \cdot \tfrac{499}{500} = \tfrac{398 - 998}{500} = -\tfrac{600}{500} = -\$1.20$.
  \pause
  \item \textbf{Reading it:} every ticket costs the buyer \$1.20 on average, and the school keeps \$600 of the \$1000 raised. No single buyer ever loses \$1.20 -- they lose \$2 or win \$398 -- yet the school can budget on it exactly.
\end{itemize}
\end{frame}

\begin{frame}{Expectation Is Where the Sample Mean Lands}
\begin{columns}[T]
\column{0.50\textwidth}
\small
\begin{itemize}
  \item Lecture 1's $\bar{x}$ was computed \emph{from data we had}. $E[X]$ is a property of the \emph{process}, fixed before any data exists.
  \pause
  \item Roll a die again and again, keeping the running $\bar{x}$. Three independent runs wander at first, then settle on $E[X] = 3.5$.
  \pause
  \item After $10$ rolls the three runs read $4.10$, $3.60$, $3.90$; after $2000$ they read $3.54$, $3.51$, $3.49$.
  \pause
  \item \textbf{Why this matters:} it is why measuring more data helps at all. Lecture 5 makes the rate precise -- the law of large numbers and the CLT.
\end{itemize}
\column{0.48\textwidth}
\centering
\includegraphics[width=\linewidth]{figures/discrete_rv/discrete_rv_running_mean.pdf}
\end{columns}
\end{frame}

% ─── Variance ───────────────────────────────────────────────────────────────
\begin{frame}{Why Do We Need the Variance of a Random Variable?}
\begin{rldefinition}
The \textbf{variance} is the expected squared distance from the mean, written $\sigma^2$:
\[
\mathrm{Var}(X) = E\!\left[(X - \mu)^2\right] = E[X^2] - \mu^2 ,
\qquad \sigma = \sqrt{\mathrm{Var}(X)} .
\]
The second form is the one to compute with: square the values, average them, subtract the square of the mean.
\end{rldefinition}
\pause
\begin{itemize}
  \item \textbf{Coffee shop:} $E[X^2] = 0^2(0.7) + 1^2(0.2) + 2^2(0.1) = 0.6$, so $\mathrm{Var}(X) = 0.6 - 0.4^2 = 0.44$ and $\sigma \approx 0.66$.
  \pause
  \item Same idea as Lecture 1's $s^2$, and same reason: the mean alone cannot tell a steady process from an erratic one.
\end{itemize}
\end{frame}

\begin{frame}{Worked Example: Two Raffles, Same Expectation}
\begin{itemize}
  \item \textbf{Raffle A} (from the previous slide): one \$400 prize among $500$ tickets, $E[X] = -\$1.20$.
  \pause
  \item \textbf{Raffle B}: same $500$ tickets at \$2, but \emph{fifty} \$8 prizes. Net gain is $+\$6$ with probability $\tfrac{50}{500} = 0.1$, else $-\$2$:
  $E[X] = 6(0.1) - 2(0.9) = -\$1.20$ -- \emph{identical}.
  \pause
  \item But $\mathrm{Var}_A = 320.8 - 1.44 = 319.36$, giving $\sigma_A \approx \$17.87$, while $\mathrm{Var}_B = 7.2 - 1.44 = 5.76$, giving $\sigma_B = \$2.40$.
  \pause
  \item \textbf{Why:} the expectation says what the school raises; the variance says what the \emph{buyer} experiences. Raffle A is a lottery, Raffle B is a mild discount -- and no expectation can tell them apart.
\end{itemize}
\end{frame}

\begin{frame}{Shifting and Scaling}
\begin{rldefinition}
For constants $a$ and $b$:
\[
E[aX + b] = a\,E[X] + b ,
\qquad
\mathrm{Var}(aX + b) = a^2\,\mathrm{Var}(X) .
\]
Adding $b$ moves the whole distribution without stretching it, so the spread is untouched; the $a^2$ appears because variance is measured in squared units.
\end{rldefinition}
\pause
\begin{itemize}
  \item \textbf{Coffee shop:} a pastry costs \$2, so pastry revenue is $R = 2X$ with $E[R] = \$0.80$ and $\sigma_R = 2(0.66) = \$1.33$.
  \pause
  \item Every order also includes a \$3 coffee, so the bill is $T = 2X + 3$: $E[T] = \$3.80$, and $\mathrm{Var}(T) = 1.76$ -- exactly $\mathrm{Var}(R)$, because a fixed \$3 adds no uncertainty.
\end{itemize}
\end{frame}

% ═══════════════════════════════════════════════════════════════════════════
